%%% ==========================================================================
%%%  第一章：实数系与确界原理
%%%  本章同时用作模板样章，演示讲义可用的全部排版构件。
%%% ==========================================================================
\chapter{实数系与确界原理}
\label{ch:reals}

\begin{keypoint}
  \begin{itemize}
    \item 有理数域对四则运算封闭，但存在「空隙」，不足以支撑极限理论。
    \item 实数系由有序域公理加上确界公理刻画，确界公理即实数的完备性。
    \item 确界原理是后续单调有界定理、闭区间套定理等一切收敛性结论的根基。
  \end{itemize}
\end{keypoint}

\section{有理数的不足}
\label{sec:rational-gap}

数学分析的全部内容都建立在实数系之上。要说明为何需要实数而非有理数，
先看一个古典结论。

\begin{proposition}
  \label{prop:sqrt2-irrational}
  不存在有理数 \(r\) 使得 \(r^{2} = 2\)。
\end{proposition}

\begin{proof}
  反设存在既约分数 \(r = p/q\)（\(p, q \in \Z\)，\(q \neq 0\)，\(\gcd(p,q)=1\)）
  使得 \(p^{2} = 2q^{2}\)。则 \(p^{2}\) 为偶数，故 \(p\) 为偶数，
  记 \(p = 2k\)。代入得 \(4k^{2} = 2q^{2}\)，即 \(q^{2} = 2k^{2}\)，
  于是 \(q\) 亦为偶数。这与 \(\gcd(p,q) = 1\) 矛盾。
\end{proof}

\Cref{prop:sqrt2-irrational} 表明，有理数集 \(\Q\) 在数轴上留有「空隙」。
集合 \(\setst{x \in \Q}{x > 0,\ x^{2} < 2}\) 在 \(\Q\) 内有上界却没有最小上界，
而分析学中的极限、连续、导数、积分等概念无一不依赖于这种最小上界的存在性。

\section{确界与确界原理}
\label{sec:supremum}

\begin{definition}[上界与上确界]
  \label{def:supremum}
  设 \(S \subseteq \R\) 非空。若存在 \(M \in \R\) 使得对一切 \(x \in S\)
  都有 \(x \le M\)，则称 \(M\) 为 \(S\) 的一个\term{上界}{upper bound}，
  并称 \(S\) 有上界。若 \(S\) 的上界之中存在最小者 \(\beta\)，
  则称 \(\beta\) 为 \(S\) 的\term{上确界}{supremum}，记作 \(\beta = \sup S\)。
\end{definition}

\begin{remark}
  按 \cref{def:supremum}，\(\beta = \sup S\) 等价于下述两条同时成立：
  \begin{enumerate}[label=(\roman*)]
    \item 对一切 \(x \in S\)，有 \(x \le \beta\)；
    \item 对任意 \(\varepsilon > 0\)，存在 \(x_{0} \in S\) 使 \(x_{0} > \beta - \varepsilon\)。
  \end{enumerate}
  第二条是实际证明中最常用的刻画：它把「最小」这一比较性陈述，
  转化为可以逐点检验的存在性陈述。
\end{remark}

\begin{assumption}[确界公理]
  \label{ax:completeness}
  \(\R\) 中任一非空有上界的子集必有上确界。
\end{assumption}

\term{确界公理}{completeness axiom}是实数区别于有理数的唯一实质性质：
\(\R\) 与 \(\Q\) 同为有序域，差别仅在于前者满足 \cref{ax:completeness}。
由此可导出全部收敛性定理。

\begin{theorem}[单调有界定理]
  \label{thm:monotone-convergence}
  设数列 \(\{a_{n}\}\) 单调递增且有上界，则 \(\{a_{n}\}\) 收敛，
  且 \(\lim_{n \to \infty} a_{n} = \sup_{n} a_{n}\)。
\end{theorem}

\begin{proof}
  记 \(S = \setst{a_{n}}{n \in \N}\)。由假设 \(S\) 非空且有上界，
  据 \cref{ax:completeness}，\(\beta = \sup S\) 存在。

  任取 \(\varepsilon > 0\)。由上确界的第二条刻画，存在 \(N\) 使
  \(a_{N} > \beta - \varepsilon\)。又 \(\{a_{n}\}\) 单调递增，
  故当 \(n > N\) 时
  \begin{equation}
    \label{eq:monotone-squeeze}
    \beta - \varepsilon < a_{N} \le a_{n} \le \beta < \beta + \varepsilon,
  \end{equation}
  即 \(\abs{a_{n} - \beta} < \varepsilon\)。由 \(\varepsilon\) 的任意性，
  \(a_{n} \to \beta\)。
\end{proof}

\begin{example}
  \label{ex:e-limit}
  证明数列 \(a_{n} = \left(1 + \dfrac{1}{n}\right)^{n}\) 收敛。
  \begin{solution}
    由二项式定理，
    \[
      a_{n} = \sum_{k=0}^{n} \binom{n}{k} \frac{1}{n^{k}}
            = \sum_{k=0}^{n} \frac{1}{k!}
              \prod_{j=1}^{k-1}\left(1 - \frac{j}{n}\right).
    \]
    把 \(n\) 换成 \(n+1\)，右端每一项中的因子 \(1 - j/n\) 均增大，
    且求和多出一项非负项，故 \(a_{n+1} > a_{n}\)，数列单调递增。

    又因 \(\prod_{j=1}^{k-1}(1 - j/n) \le 1\) 且 \(k! \ge 2^{k-1}\)，
    \[
      a_{n} \le \sum_{k=0}^{n} \frac{1}{k!}
            \le 1 + \sum_{k=1}^{n} \frac{1}{2^{k-1}} < 3,
    \]
    数列有上界。由 \cref{thm:monotone-convergence} 即得收敛。
    该极限记作 \(\eu\)。
  \end{solution}
\end{example}

关于实数系的公理化构造与戴德金分割的等价性，可参见 \textcite{rudin1976} 第一章。

\section*{习题}
\addcontentsline{toc}{section}{习题}

\begin{exercise}
  \label{exr:inf-def}
  仿照 \cref{def:supremum} 写出下界与下确界 \(\inf S\) 的定义，
  并给出与上确界第二条刻画相对应的 \(\varepsilon\) 形式。
  \answerspace[3cm]
  \begin{solution}
    设 \(S \subseteq \R\) 非空。若存在 \(m \in \R\) 使对一切 \(x \in S\)
    有 \(x \ge m\)，则称 \(m\) 为 \(S\) 的下界。\(S\) 的下界之中的最大者
    \(\alpha\) 称为 \(S\) 的下确界，记 \(\alpha = \inf S\)。

    \(\varepsilon\) 形式：\(\alpha = \inf S\) 当且仅当
    （i）对一切 \(x \in S\) 有 \(x \ge \alpha\)；
    （ii）对任意 \(\varepsilon > 0\)，存在 \(x_{0} \in S\) 使
    \(x_{0} < \alpha + \varepsilon\)。
  \end{solution}
\end{exercise}

\begin{exercise}
  \label{exr:sup-sum}
  设 \(A, B \subseteq \R\) 非空且均有上界，令
  \(A + B = \setst{a + b}{a \in A,\ b \in B}\)。
  证明：\(\sup(A + B) = \sup A + \sup B\)。
  \answerspace[4.5cm]
  \begin{solution}
    记 \(\alpha = \sup A\)，\(\beta = \sup B\)。

    先证 \(\le\)。任取 \(a \in A\)、\(b \in B\)，有
    \(a + b \le \alpha + \beta\)，故 \(\alpha + \beta\) 是 \(A + B\) 的上界，
    从而 \(\sup(A+B) \le \alpha + \beta\)。

    再证 \(\ge\)。任取 \(\varepsilon > 0\)。存在 \(a_{0} \in A\) 使
    \(a_{0} > \alpha - \varepsilon/2\)，存在 \(b_{0} \in B\) 使
    \(b_{0} > \beta - \varepsilon/2\)。于是
    \(a_{0} + b_{0} > \alpha + \beta - \varepsilon\)，
    故 \(\sup(A+B) > \alpha + \beta - \varepsilon\)。
    由 \(\varepsilon\) 的任意性得 \(\sup(A+B) \ge \alpha + \beta\)。

    两式合并即得结论。
  \end{solution}
\end{exercise}
